
\documentclass{article}
\usepackage{colortbl}
\usepackage{makecell}
\usepackage{multirow}
\usepackage{supertabular}

\begin{document}

\newcounter{utterance}

\centering \large Interaction Transcript for game `function\_detective', experiment `two\_numbers\_passive\_labeled\_pairs', episode 243 with Qwen3.6{-}35B{-}A3B{-}FP8{-}without{-}reasoning.
\vspace{24pt}

{ \footnotesize  \setcounter{utterance}{1}
\setlength{\tabcolsep}{0pt}
\begin{supertabular}{c@{$\;$}|p{.15\linewidth}@{}p{.15\linewidth}p{.15\linewidth}p{.15\linewidth}p{.15\linewidth}p{.15\linewidth}}
    \# & \multicolumn{2}{c}{Player} && \multicolumn{2}{c}{Game Master} \\
    \hline

    \theutterance \stepcounter{utterance}  
    & & & \multicolumn{4}{p{0.6\linewidth}}{
        \cellcolor[rgb]{0.9,0.9,0.9}{
            \makecell[{{p{\linewidth}}}]{
                \texttt{\tiny{[P1$\langle$GM]}}
                \texttt{I have a hidden Python function for you to figure out.} \\
\\ 
\texttt{Your job is to infer what it does and then submit a correct Python implementation.} \\
\\ 
\texttt{Function signature:} \\
\texttt{(a: int, b: int) {-}> int} \\
\\ 
\texttt{You have 15 turns total.} \\
\\ 
\texttt{I reveal candidate input{-}output pairs one at a time and tell you whether each one is correct.} \\
\\ 
\texttt{You must make your best SOLVE attempt no later than the second{-}to{-}last turn.} \\
\texttt{Do not wait until the very end.} \\
\\ 
\texttt{You can use these actions:} \\
\texttt{{-} JSON format described at the end of the prompt} \\
\texttt{{-} SOLVE: ```python ... ```} \\
\\ 
\texttt{Do not add import statements in your guess.} \\
\texttt{Anything you might need is already available, so just write the function body in hypothesis and return the answer.} \\
\\ 
\texttt{Preloaded examples:} \\
\texttt{1. (((99, 1000), 9), True)} \\
\\ 
\texttt{Scoring:} \\
\texttt{{-} correctness is required for any non{-}zero score} \\
\texttt{{-} in interactive modes, solving with fewer information requests is better} \\
\\ 
\texttt{Input formatting reminders:} \\
\texttt{{-} strings must be quoted} \\
\texttt{{-} booleans must be True/False} \\
\texttt{{-} use commas between multiple values} \\
\texttt{{-} do not write any extra action tags in the same message} \\
\\ 
\texttt{Format example:} \\
\texttt{You write:} \\
\texttt{\{'current\_mode': <probing/confirming/don't know; choose whichever best describes your current state in the task>, 'input\_rationale': <brief description of why you chose this hypothesis>, 'hypothesis': <either the exact string "unknown" if you have no guess yet OR your current best guess as runnable Python matching the given signature, with a one{-}line docstring stating what it does. Never write placeholder code, if you cannot commit to a guess, use "unknown". Match the given signature.} \\
\texttt{e.g.} \\
\texttt{def f(x):} \\
\texttt{	"""Returns the absolute value of x."""} \\
\texttt{	return x if x >= 0 else {-}x} \\
\\ 
\texttt{Use "unknown" if you have no guess.>, 'confidence\_score': <score between 1{-}100. 100 means highest confidence in your hypothesis>\}} \\
\\ 
\texttt{Output only one valid json per turn.} \\
            }
        }
    }
     \\ \\

    \theutterance \stepcounter{utterance}  
    & \multicolumn{4}{p{0.6\linewidth}}{
        \cellcolor[rgb]{0.561,0.737,0.561}{
            \makecell[{{p{\linewidth}}}]{
                \texttt{\tiny{[P1$\rangle$GM]}}
                \texttt{\{'current\_mode': 'probing', 'input\_rationale': 'I have one preloaded example: f(99, 1000) = True. This suggests the function might return a boolean, but the signature says it returns an int. In Python, True is equivalent to 1 and False to 0 when returned as an int. Let me probe with simple inputs to understand the logic. I'll try small integers to see if there is a pattern like comparison, modulo, or arithmetic operation.', 'hypothesis': 'unknown', 'confidence\_score': 10\}} \\
            }
        }
    }
    & & \\ \\

    \theutterance \stepcounter{utterance}  
    & & & \multicolumn{4}{p{0.6\linewidth}}{
        \cellcolor[rgb]{0.9,0.9,0.9}{
            \makecell[{{p{\linewidth}}}]{
                \texttt{\tiny{[P1$\langle$GM]}}
                \texttt{Game Violation: ParseError: Response is not valid JSON.} \\
\texttt{Output ONLY:} \\
\texttt{One JSON with confidence\_score, current\_mode, input\_rationale, and hypothesis keys ...} \\
\texttt{OR} \\
\texttt{SOLVE: ```python ... ```} \\
            }
        }
    }
     \\ \\

    \theutterance \stepcounter{utterance}  
    & \multicolumn{4}{p{0.6\linewidth}}{
        \cellcolor[rgb]{0.561,0.737,0.561}{
            \makecell[{{p{\linewidth}}}]{
                \texttt{\tiny{[P1$\rangle$GM]}}
                \texttt{\{} \\
\texttt{"current\_mode": "probing",} \\
\texttt{"input\_rationale": "I have one preloaded example: f(99, 1000) = True. Since the return type is int, True likely means 1. I need to probe with more inputs to understand the relationship between a and b. I will try simple cases like (0, 0), (1, 1), (0, 1), (1, 0) to see if it's a comparison or logical operation.",} \\
\texttt{"hypothesis": "unknown",} \\
\texttt{"confidence\_score": 10} \\
\texttt{\}} \\
            }
        }
    }
    & & \\ \\

    \theutterance \stepcounter{utterance}  
    & & & \multicolumn{4}{p{0.6\linewidth}}{
        \cellcolor[rgb]{0.9,0.9,0.9}{
            \makecell[{{p{\linewidth}}}]{
                \texttt{\tiny{[P1$\langle$GM]}}
                \texttt{OUTPUT: ((({-}99, 0), 19), False)} \\
            }
        }
    }
     \\ \\

    \theutterance \stepcounter{utterance}  
    & \multicolumn{4}{p{0.6\linewidth}}{
        \cellcolor[rgb]{0.561,0.737,0.561}{
            \makecell[{{p{\linewidth}}}]{
                \texttt{\tiny{[P1$\rangle$GM]}}
                \texttt{\{} \\
\texttt{"current\_mode": "probing",} \\
\texttt{"input\_rationale": "I have two data points: f(99, 1000) = True (1) and f({-}99, 0) = False (0). Let's analyze. 99 and 1000 are both positive. {-}99 is negative, 0 is zero. Maybe it checks if both are positive? Or if a < b? 99 < 1000 is True. {-}99 < 0 is True, but result is False. So it's not just a < b. Maybe it checks if both are non{-}negative? 99 >= 0 and 1000 >= 0 is True. {-}99 >= 0 is False. This fits. Let's test with another case where one is negative and one is positive, or both negative.",} \\
\texttt{"hypothesis": "unknown",} \\
\texttt{"confidence\_score": 30} \\
\texttt{\}} \\
            }
        }
    }
    & & \\ \\

    \theutterance \stepcounter{utterance}  
    & & & \multicolumn{4}{p{0.6\linewidth}}{
        \cellcolor[rgb]{0.9,0.9,0.9}{
            \makecell[{{p{\linewidth}}}]{
                \texttt{\tiny{[P1$\langle$GM]}}
                \texttt{OUTPUT: (((2, {-}2), 4), True)} \\
            }
        }
    }
     \\ \\

    \theutterance \stepcounter{utterance}  
    & \multicolumn{4}{p{0.6\linewidth}}{
        \cellcolor[rgb]{0.561,0.737,0.561}{
            \makecell[{{p{\linewidth}}}]{
                \texttt{\tiny{[P1$\rangle$GM]}}
                \texttt{\{} \\
\texttt{"current\_mode": "probing",} \\
\texttt{"input\_rationale": "I have three data points: f(99, 1000) = True, f({-}99, 0) = False, f(2, {-}2) = True. The hypothesis that it checks if both are non{-}negative is incorrect because f(2, {-}2) is True but {-}2 is negative. Let's look at the values. 99, 1000 {-}> True. {-}99, 0 {-}> False. 2, {-}2 {-}> True. Maybe it's about the sum? 99+1000=1099 (positive). {-}99+0={-}99 (negative). 2+({-}2)=0 (zero). If it checks if sum > 0, then 0 would be False. Let's test if sum >= 0. Or maybe it's about absolute values? |99| < |1000|? |{-}99| > |0|? |2| == |{-}2|? This seems inconsistent. Let's try to find a pattern with modulo or bitwise operations. Or perhaps it's checking if a and b have the same sign? 99, 1000 same sign (pos) {-}> True. {-}99, 0 different signs? 0 is neither. 2, {-}2 different signs {-}> True.} \\
            }
        }
    }
    & & \\ \\

    \theutterance \stepcounter{utterance}  
    & & & \multicolumn{4}{p{0.6\linewidth}}{
        \cellcolor[rgb]{0.9,0.9,0.9}{
            \makecell[{{p{\linewidth}}}]{
                \texttt{\tiny{[P1$\langle$GM]}}
                \texttt{Game Violation: ParseError: Response is not valid JSON.} \\
\texttt{Output ONLY:} \\
\texttt{One JSON with confidence\_score, current\_mode, input\_rationale, and hypothesis keys ...} \\
\texttt{OR} \\
\texttt{SOLVE: ```python ... ```} \\
            }
        }
    }
     \\ \\

    \theutterance \stepcounter{utterance}  
    & \multicolumn{4}{p{0.6\linewidth}}{
        \cellcolor[rgb]{0.561,0.737,0.561}{
            \makecell[{{p{\linewidth}}}]{
                \texttt{\tiny{[P1$\rangle$GM]}}
                \texttt{\{} \\
\texttt{"current\_mode": "probing",} \\
\texttt{"input\_rationale": "I have three data points: f(99, 1000) = True, f({-}99, 0) = False, f(2, {-}2) = True. The hypothesis that it checks if both are non{-}negative is incorrect because f(2, {-}2) is True but {-}2 is negative. Let's look at the values. 99, 1000 {-}> True. {-}99, 0 {-}> False. 2, {-}2 {-}> True. Maybe it's about the sum? 99+1000=1099 (positive). {-}99+0={-}99 (negative). 2+({-}2)=0 (zero). If it checks if sum > 0, then 0 would be False. Let's test if sum >= 0. Or maybe it's about absolute values? |99| < |1000|? |{-}99| > |0|? |2| == |{-}2|? This seems inconsistent. Let's try to find a pattern with modulo or bitwise operations. Or perhaps it's checking if a and b have the same sign? 99, 1000 same sign (pos) {-}> True. {-}99, 0 different signs? 0 is neither. 2, {-}2 different signs {-}> True.",} \\
            }
        }
    }
    & & \\ \\

    \theutterance \stepcounter{utterance}  
    & & & \multicolumn{4}{p{0.6\linewidth}}{
        \cellcolor[rgb]{0.9,0.9,0.9}{
            \makecell[{{p{\linewidth}}}]{
                \texttt{\tiny{[P1$\langle$GM]}}
                \texttt{Game Violation: ParseError: Response is not valid JSON.} \\
\texttt{Output ONLY:} \\
\texttt{One JSON with confidence\_score, current\_mode, input\_rationale, and hypothesis keys ...} \\
\texttt{OR} \\
\texttt{SOLVE: ```python ... ```} \\
            }
        }
    }
     \\ \\

    \theutterance \stepcounter{utterance}  
    & \multicolumn{4}{p{0.6\linewidth}}{
        \cellcolor[rgb]{0.561,0.737,0.561}{
            \makecell[{{p{\linewidth}}}]{
                \texttt{\tiny{[P1$\rangle$GM]}}
                \texttt{\{} \\
\texttt{"current\_mode": "probing",} \\
\texttt{"input\_rationale": "I have three data points: f(99, 1000) = True, f({-}99, 0) = False, f(2, {-}2) = True. Let's re{-}evaluate. Maybe it's checking if a and b are coprime? gcd(99, 1000) = 1 {-}> True. gcd({-}99, 0) = 99 {-}> False. gcd(2, {-}2) = 2 {-}> True. This doesn't fit because gcd(2, {-}2) is 2, not 1. Maybe it's checking if a divides b or b divides a? 99 does not divide 1000. 1000 does not divide 99. {-}99 does not divide 0? Actually 0 is divisible by any non{-}zero integer. So {-}99 divides 0. But result is False. So that's not it. Let's look at bitwise AND. 99 \& 1000 = 0? 99 is 0b1100011, 1000 is 0b1111101000. AND is 0. {-}99 \& 0 = 0. 2 \& {-}2 = 2? In two's complement} \\
            }
        }
    }
    & & \\ \\

    \theutterance \stepcounter{utterance}  
    & & & \multicolumn{4}{p{0.6\linewidth}}{
        \cellcolor[rgb]{0.9,0.9,0.9}{
            \makecell[{{p{\linewidth}}}]{
                \texttt{\tiny{[P1$\langle$GM]}}
                \texttt{Game Violation: ParseError: Response is not valid JSON.} \\
\texttt{Output ONLY:} \\
\texttt{One JSON with confidence\_score, current\_mode, input\_rationale, and hypothesis keys ...} \\
\texttt{OR} \\
\texttt{SOLVE: ```python ... ```} \\
            }
        }
    }
     \\ \\

    \theutterance \stepcounter{utterance}  
    & \multicolumn{4}{p{0.6\linewidth}}{
        \cellcolor[rgb]{0.561,0.737,0.561}{
            \makecell[{{p{\linewidth}}}]{
                \texttt{\tiny{[P1$\rangle$GM]}}
                \texttt{\{} \\
\texttt{"current\_mode": "probing",} \\
\texttt{"input\_rationale": "I have three data points: f(99, 1000) = True, f({-}99, 0) = False, f(2, {-}2) = True. Let's look at the sum again. 99+1000=1099 (odd). {-}99+0={-}99 (odd). 2+({-}2)=0 (even). This doesn't seem to correlate with True/False directly. Let's look at the product. 99*1000=99000 (even). {-}99*0=0 (even). 2*{-}2={-}4 (even). No. Let's look at bitwise XOR. 99\^{}1000 = 939. {-}99\^{}0 = {-}99. 2\^{}{-}2 = 0. No obvious pattern. Let's try a new probe: f(1, 1). If it's 'same sign or one is zero', 1,1 is True. If it's 'sum is odd', 1+1=2 even {-}> False. If it's 'coprime', gcd(1,1)=1 {-}> True. Let's probe f(1, 1) and f(0, 0).",} \\
\texttt{"hypothesis": "unknown",} \\
\texttt{"} \\
            }
        }
    }
    & & \\ \\

    \theutterance \stepcounter{utterance}  
    & & & \multicolumn{4}{p{0.6\linewidth}}{
        \cellcolor[rgb]{0.9,0.9,0.9}{
            \makecell[{{p{\linewidth}}}]{
                \texttt{\tiny{[P1$\langle$GM]}}
                \texttt{Game Violation: ParseError: Response is not valid JSON.} \\
\texttt{Output ONLY:} \\
\texttt{One JSON with confidence\_score, current\_mode, input\_rationale, and hypothesis keys ...} \\
\texttt{OR} \\
\texttt{SOLVE: ```python ... ```} \\
            }
        }
    }
     \\ \\

    \theutterance \stepcounter{utterance}  
    & \multicolumn{4}{p{0.6\linewidth}}{
        \cellcolor[rgb]{0.561,0.737,0.561}{
            \makecell[{{p{\linewidth}}}]{
                \texttt{\tiny{[P1$\rangle$GM]}}
                \texttt{\{} \\
\texttt{"current\_mode": "probing",} \\
\texttt{"input\_rationale": "I have three data points: f(99, 1000) = True, f({-}99, 0) = False, f(2, {-}2) = True. I need more data to distinguish between hypotheses like 'sum is odd', 'coprime', 'bitwise AND is 0', etc. I will probe with (1, 1) and (0, 0) to check edge cases and parity.",} \\
\texttt{"hypothesis": "unknown",} \\
\texttt{"confidence\_score": 20} \\
\texttt{\}} \\
            }
        }
    }
    & & \\ \\

    \theutterance \stepcounter{utterance}  
    & & & \multicolumn{4}{p{0.6\linewidth}}{
        \cellcolor[rgb]{0.9,0.9,0.9}{
            \makecell[{{p{\linewidth}}}]{
                \texttt{\tiny{[P1$\langle$GM]}}
                \texttt{OUTPUT: ((({-}1000, {-}1000), {-}1), False)} \\
            }
        }
    }
     \\ \\

    \theutterance \stepcounter{utterance}  
    & \multicolumn{4}{p{0.6\linewidth}}{
        \cellcolor[rgb]{0.561,0.737,0.561}{
            \makecell[{{p{\linewidth}}}]{
                \texttt{\tiny{[P1$\rangle$GM]}}
                \texttt{\{} \\
\texttt{"current\_mode": "probing",} \\
\texttt{"input\_rationale": "I have four data points: f(99, 1000) = True, f({-}99, 0) = False, f(2, {-}2) = True, f({-}1000, {-}1000) = False. Let's re{-}evaluate. f(2, {-}2) is True. f({-}1000, {-}1000) is False. Both pairs have equal absolute values. The difference is signs. 2 and {-}2 have different signs. {-}1000 and {-}1000 have same signs. f(99, 1000) different signs? No, both positive. f({-}99, 0) different signs? {-}99 neg, 0 neutral. Let's look at bitwise AND again. 99 \& 1000 = 0. {-}99 \& 0 = 0. 2 \& {-}2 = 2 (non{-}zero). {-}1000 \& {-}1000 = {-}1000 (non{-}zero). This doesn't fit 'AND is 0'. Let's look at bitwise OR. 99 | 1000 = 1007. {-}99 | 0 = {-}99. 2 | {-}2 = {-}2. {-}1} \\
            }
        }
    }
    & & \\ \\

    \theutterance \stepcounter{utterance}  
    & & & \multicolumn{4}{p{0.6\linewidth}}{
        \cellcolor[rgb]{0.9,0.9,0.9}{
            \makecell[{{p{\linewidth}}}]{
                \texttt{\tiny{[P1$\langle$GM]}}
                \texttt{Game Violation: ParseError: Response is not valid JSON.} \\
\texttt{Output ONLY:} \\
\texttt{One JSON with confidence\_score, current\_mode, input\_rationale, and hypothesis keys ...} \\
\texttt{OR} \\
\texttt{SOLVE: ```python ... ```} \\
            }
        }
    }
     \\ \\

    \theutterance \stepcounter{utterance}  
    & \multicolumn{4}{p{0.6\linewidth}}{
        \cellcolor[rgb]{0.561,0.737,0.561}{
            \makecell[{{p{\linewidth}}}]{
                \texttt{\tiny{[P1$\rangle$GM]}}
                \texttt{\{} \\
\texttt{"current\_mode": "probing",} \\
\texttt{"input\_rationale": "I have four data points: f(99, 1000) = True, f({-}99, 0) = False, f(2, {-}2) = True, f({-}1000, {-}1000) = False. Let's look at the bitwise AND of a and b. 99 \& 1000 = 0. {-}99 \& 0 = 0. 2 \& {-}2 = 2 (non{-}zero). {-}1000 \& {-}1000 = {-}1000 (non{-}zero). This doesn't fit a simple 'AND is 0' rule. Let's look at the sum. 99+1000=1099 (odd). {-}99+0={-}99 (odd). 2{-}2=0 (even). {-}1000{-}1000={-}2000 (even). True for odd sums, False for even sums? Let's check: 1099 is odd {-}> True. {-}99 is odd {-}> False. This contradicts. Wait, {-}99 is odd. So sum parity doesn't work. Let's look at bitwise XOR. 99\^{}1000 = 939. {-}99\^{}0 = {-}99.} \\
            }
        }
    }
    & & \\ \\

    \theutterance \stepcounter{utterance}  
    & & & \multicolumn{4}{p{0.6\linewidth}}{
        \cellcolor[rgb]{0.9,0.9,0.9}{
            \makecell[{{p{\linewidth}}}]{
                \texttt{\tiny{[P1$\langle$GM]}}
                \texttt{Game Violation: ParseError: Response is not valid JSON.} \\
\texttt{Output ONLY:} \\
\texttt{One JSON with confidence\_score, current\_mode, input\_rationale, and hypothesis keys ...} \\
\texttt{OR} \\
\texttt{SOLVE: ```python ... ```} \\
            }
        }
    }
     \\ \\

    \theutterance \stepcounter{utterance}  
    & \multicolumn{4}{p{0.6\linewidth}}{
        \cellcolor[rgb]{0.561,0.737,0.561}{
            \makecell[{{p{\linewidth}}}]{
                \texttt{\tiny{[P1$\rangle$GM]}}
                \texttt{\{} \\
\texttt{"current\_mode": "probing",} \\
\texttt{"input\_rationale": "I have four data points: f(99, 1000) = True, f({-}99, 0) = False, f(2, {-}2) = True, f({-}1000, {-}1000) = False. Let's look at the bitwise AND of a and b. 99 \& 1000 = 0. {-}99 \& 0 = 0. 2 \& {-}2 = 2 (non{-}zero). {-}1000 \& {-}1000 = {-}1000 (non{-}zero). This doesn't fit a simple 'AND is 0' rule. Let's look at the sum. 99+1000=1099 (odd). {-}99+0={-}99 (odd). 2{-}2=0 (even). {-}1000{-}1000={-}2000 (even). True for odd sums, False for even sums? Let's check: 1099 is odd {-}> True. {-}99 is odd {-}> False. This contradicts. Wait, {-}99 is odd. So sum parity doesn't work. Let's look at bitwise XOR. 99\^{}1000 = 939. {-}99\^{}0 = {-}99.} \\
            }
        }
    }
    & & \\ \\

    \theutterance \stepcounter{utterance}  
    & & & \multicolumn{4}{p{0.6\linewidth}}{
        \cellcolor[rgb]{0.9,0.9,0.9}{
            \makecell[{{p{\linewidth}}}]{
                \texttt{\tiny{[P1$\langle$GM]}}
                \texttt{Game Violation: ParseError: Response is not valid JSON.} \\
\texttt{Output ONLY:} \\
\texttt{One JSON with confidence\_score, current\_mode, input\_rationale, and hypothesis keys ...} \\
\texttt{OR} \\
\texttt{SOLVE: ```python ... ```} \\
            }
        }
    }
     \\ \\

    \theutterance \stepcounter{utterance}  
    & \multicolumn{4}{p{0.6\linewidth}}{
        \cellcolor[rgb]{0.561,0.737,0.561}{
            \makecell[{{p{\linewidth}}}]{
                \texttt{\tiny{[P1$\rangle$GM]}}
                \texttt{\{} \\
\texttt{"current\_mode": "probing",} \\
\texttt{"input\_rationale": "I have four data points: f(99, 1000) = True, f({-}99, 0) = False, f(2, {-}2) = True, f({-}1000, {-}1000) = False. Let's look at the bitwise AND of a and b. 99 \& 1000 = 0. {-}99 \& 0 = 0. 2 \& {-}2 = 2 (non{-}zero). {-}1000 \& {-}1000 = {-}1000 (non{-}zero). This doesn't fit a simple 'AND is 0' rule. Let's look at the sum. 99+1000=1099 (odd). {-}99+0={-}99 (odd). 2{-}2=0 (even). {-}1000{-}1000={-}2000 (even). True for odd sums, False for even sums? Let's check: 1099 is odd {-}> True. {-}99 is odd {-}> False. This contradicts. Wait, {-}99 is odd. So sum parity doesn't work. Let's look at bitwise XOR. 99\^{}1000 = 939. {-}99\^{}0 = {-}99.} \\
            }
        }
    }
    & & \\ \\

    \theutterance \stepcounter{utterance}  
    & & & \multicolumn{4}{p{0.6\linewidth}}{
        \cellcolor[rgb]{0.9,0.9,0.9}{
            \makecell[{{p{\linewidth}}}]{
                \texttt{\tiny{[P1$\langle$GM]}}
                \texttt{Game Violation: ParseError: Response is not valid JSON.} \\
\texttt{Output ONLY:} \\
\texttt{One JSON with confidence\_score, current\_mode, input\_rationale, and hypothesis keys ...} \\
\texttt{OR} \\
\texttt{SOLVE: ```python ... ```} \\
            }
        }
    }
     \\ \\

    \theutterance \stepcounter{utterance}  
    & \multicolumn{4}{p{0.6\linewidth}}{
        \cellcolor[rgb]{0.561,0.737,0.561}{
            \makecell[{{p{\linewidth}}}]{
                \texttt{\tiny{[P1$\rangle$GM]}}
                \texttt{\{} \\
\texttt{"current\_mode": "probing",} \\
\texttt{"input\_rationale": "I have four data points: f(99, 1000) = True, f({-}99, 0) = False, f(2, {-}2) = True, f({-}1000, {-}1000) = False. Let's look at the bitwise AND of a and b. 99 \& 1000 = 0. {-}99 \& 0 = 0. 2 \& {-}2 = 2 (non{-}zero). {-}1000 \& {-}1000 = {-}1000 (non{-}zero). This doesn't fit a simple 'AND is 0' rule. Let's look at the sum. 99+1000=1099 (odd). {-}99+0={-}99 (odd). 2{-}2=0 (even). {-}1000{-}1000={-}2000 (even). True for odd sums, False for even sums? Let's check: 1099 is odd {-}> True. {-}99 is odd {-}> False. This contradicts. Wait, {-}99 is odd. So sum parity doesn't work. Let's look at bitwise XOR. 99\^{}1000 = 939. {-}99\^{}0 = {-}99.",} \\
            }
        }
    }
    & & \\ \\

    \theutterance \stepcounter{utterance}  
    & & & \multicolumn{4}{p{0.6\linewidth}}{
        \cellcolor[rgb]{0.9,0.9,0.9}{
            \makecell[{{p{\linewidth}}}]{
                \texttt{\tiny{[P1$\langle$GM]}}
                \texttt{Game Violation: ParseError: Response is not valid JSON.} \\
\texttt{Output ONLY:} \\
\texttt{One JSON with confidence\_score, current\_mode, input\_rationale, and hypothesis keys ...} \\
\texttt{OR} \\
\texttt{SOLVE: ```python ... ```} \\
            }
        }
    }
     \\ \\

    \theutterance \stepcounter{utterance}  
    & \multicolumn{4}{p{0.6\linewidth}}{
        \cellcolor[rgb]{0.561,0.737,0.561}{
            \makecell[{{p{\linewidth}}}]{
                \texttt{\tiny{[P1$\rangle$GM]}}
                \texttt{\{} \\
\texttt{"current\_mode": "probing",} \\
\texttt{"input\_rationale": "I have four data points: f(99, 1000) = True, f({-}99, 0) = False, f(2, {-}2) = True, f({-}1000, {-}1000) = False. Let's look at the bitwise AND of a and b. 99 \& 1000 = 0. {-}99 \& 0 = 0. 2 \& {-}2 = 2 (non{-}zero). {-}1000 \& {-}1000 = {-}1000 (non{-}zero). This doesn't fit a simple 'AND is 0' rule. Let's look at the sum. 99+1000=1099 (odd). {-}99+0={-}99 (odd). 2{-}2=0 (even). {-}1000{-}1000={-}2000 (even). True for odd sums, False for even sums? Let's check: 1099 is odd {-}> True. {-}99 is odd {-}> False. This contradicts. Wait, {-}99 is odd. So sum parity doesn't work. Let's look at bitwise XOR. 99\^{}1000 = 939. {-}99\^{}0 = {-}99.",} \\
            }
        }
    }
    & & \\ \\

    \theutterance \stepcounter{utterance}  
    & & & \multicolumn{4}{p{0.6\linewidth}}{
        \cellcolor[rgb]{0.9,0.9,0.9}{
            \makecell[{{p{\linewidth}}}]{
                \texttt{\tiny{[P1$\langle$GM]}}
                \texttt{Game Violation: ParseError: Response is not valid JSON.} \\
\texttt{Output ONLY:} \\
\texttt{One JSON with confidence\_score, current\_mode, input\_rationale, and hypothesis keys ...} \\
\texttt{OR} \\
\texttt{SOLVE: ```python ... ```} \\
            }
        }
    }
     \\ \\

    \theutterance \stepcounter{utterance}  
    & \multicolumn{4}{p{0.6\linewidth}}{
        \cellcolor[rgb]{0.561,0.737,0.561}{
            \makecell[{{p{\linewidth}}}]{
                \texttt{\tiny{[P1$\rangle$GM]}}
                \texttt{\{} \\
\texttt{"current\_mode": "probing",} \\
\texttt{"input\_rationale": "I have four data points: f(99, 1000) = True, f({-}99, 0) = False, f(2, {-}2) = True, f({-}1000, {-}1000) = False. Let's look at the bitwise AND of a and b. 99 \& 1000 = 0. {-}99 \& 0 = 0. 2 \& {-}2 = 2 (non{-}zero). {-}1000 \& {-}1000 = {-}1000 (non{-}zero). This doesn't fit a simple 'AND is 0' rule. Let's look at the sum. 99+1000=1099 (odd). {-}99+0={-}99 (odd). 2{-}2=0 (even). {-}1000{-}1000={-}2000 (even). True for odd sums, False for even sums? Let's check: 1099 is odd {-}> True. {-}99 is odd {-}> False. This contradicts. Wait, {-}99 is odd. So sum parity doesn't work. Let's look at bitwise XOR. 99\^{}1000 = 939. {-}99\^{}0 = {-}99.",} \\
            }
        }
    }
    & & \\ \\

    \theutterance \stepcounter{utterance}  
    & & & \multicolumn{2}{p{0.3\linewidth}}{
        \cellcolor[rgb]{0.9,0.9,0.9}{
            \makecell[{{p{\linewidth}}}]{
                \texttt{\tiny{[GM$|$GM]}}
                \texttt{Game Verdict: LOSS (turn limit reached)} \\
            }
        }
    }
    & & \\ \\

\end{supertabular}
}

\end{document}
